\clearpage\thispagestyle{empty}
\pdfbookmark[0]{Quellen zu den Sternbedeckungen (8. März 1818)}{occult1818mar-sources}
\begin{center}\textbf{Quellen zu den Sternbedeckungen (8. März 1818)}\end{center}
\noindent Diese KI-Übertragung ins heutige Deutsch folgt den an Johann Elert Bode gesandten Beobachtungen von Carl Friedrich Gauß in den \href{https://gdz.sub.uni-goettingen.de/id/PPN236007467}{\textit{Werken}, Band~VI} (Göttingen, 1874), S.~402. Das \textit{Astronomische Jahrbuch für 1821} (Berlin, 1818), S.~\href{https://api.digitale-sammlungen.de/iiif/image/v2/bsb10538350_00238/full/1600,/0/default.jpg}{232}, nennt Gauß in Göttingen. Der Link führt zum Exemplar der Bayerischen Staatsbibliothek.
\par\smallskip\noindent Das Jahrbuch verwendet beim Januareintrag ein hochgestelltes h, einen einfachen und einen doppelten Strich und vor der abschließenden Null ein Komma. Die Werkausgabe verwendet hochgestellte u, m und s und setzt die Null unmittelbar hinter das s. Die gepaarten Beobachtungen von Gauß und Harding sowie die Zeichen für Jungfrau und Stier sind übernommen. KI-Übertragung und KI-Satz.
\clearpage
